\documentclass[11pt]{article}
\input{preamble.tex}

% placeholders for coating results that still need to be filled in
\newcommand{\TODO}[1]{\textcolor{red}{[TODO: #1]}}

\title{Optimization of wire coating}
\author{Snehal Choudhury, Eric Dickson-Peppler, Britney Sedano, Kyle Wodehouse \\(``\textit{Viscous Cycle}")}
\date{}

\begin{document}
\maketitle
\begin{abstract}
Fluoropolymers FEP 9835 and PFA 440 provide the electrical properties required for wire coating, but their melts are only mildly shear-thinning, which limits line speed before melt fracture. This study examines whether a coextruded outer layer of HT 2181, a strongly shear-thinning process aid, can raise that limit. Frequency sweeps were converted to master curves by time--temperature superposition, fit with Cross models, and used to compute line speed in a pressure-extrusion die with a 0.20~mm FEP inner layer, a total coating thickness of at most 0.25~mm, and a 330~°C limit. The HT thickness was optimized with a bounded Brent search subject to the melt-fracture and draw-down constraints. The optimum is a 0.029~mm HT layer, which allows a line speed of 0.057~m/s (3.4~m/min), 3.45 times the 0.0165~m/s of a pure 0.20~mm FEP coating at the same temperature. The optimal HT thickness is insensitive to the shifting method (0.027 to 0.029~mm), but the predicted line speed varies with the shift basis (0.054 to 0.090~m/s at 330~°C), so the speed estimate carries that uncertainty. The result also depends on extrapolating HT rheology 15~°C beyond the measured range; at 315~°C the same design gives 2.2~m/min.
\end{abstract}


\section*{Introduction}
FEP 9835 and PFA 440 are fluoropolymers that are used in wire coating due to their electrical properties that are required for the cable product (rewrite).

These fluoropolymers are only mild shear-thinning therefore we need other goodies to have successful wire coating ayyayayayay

- want to move the line faster so looking into multilayer coatings w/ non-fluoropolymer HT




\section*{Analyzing Trios data}

Master curves were assembled from frequency sweeps for FEP, PFA, and HT inside Trios. For FEP and PFA the auto-shift-finding was done on $\tan(\delta)$. For HT the auto-shift was done on storage modulus, loss modulus, and complex viscosity since the $\tan(\delta)$ shift did not make the viscosity data overlap nicely (\cref{app:tandelta_shift}). Only x-shifting was done,  $b_T$ values were assumed unity. Shifted complex viscosity master curves were fit to a Cross model\cite{cross_rheology_1965} (Eq.~\ref{eq:cross}). The fit parameters are reported in \cref{tab:cross_model_parameters}. Since the Cross model adequately fits the shifted viscosity data, a model with more parameters, namely the Carreau-Yasuda model, was not explored.

\begin{equation}
    \eta(\dot{\gamma}) = \eta_{\infty} + \frac{\eta_{0} - \eta_{\infty}}{1 + (\lambda \dot{\gamma})^{n}}
    \label{eq:cross}
\end{equation}

% NEW
The shift factors from Trios were fit to an Arrhenius form, $\ln a_T = (E_a/R)(1/T - 1/T_\text{ref})$, which gives the activation energies in \cref{tab:cross_model_parameters}. With only three temperatures per polymer these fits are coarse, but they are linear in $1/T$ ($R^2 = 0.997$ for FEP and PFA, $0.989$ for HT). HT has roughly twice the activation energy of either fluoropolymer, so its viscosity drops faster as the die gets hotter. That difference matters later, because it is why the two-layer coating gains more from heating than the single FEP layer does. The Cross fits (with $\eta_\infty = 0$) reproduce the measured complex viscosity to a median error of 2--4\% at every temperature for FEP and PFA. For HT the error is 2\% at the 300~°C reference and about 8\% at 285 and 315~°C, which reflects the imperfect overlap of the HT master curve.

Every sweep ran from high to low frequency, so the lowest-frequency points were measured last, 10--35~min into the test. In the terminal region $|\eta^*|$ should level off as $\omega$ decreases. In a few sweeps it fell instead, by up to 12\% for HT at 315~°C, which points to the sample degrading during the test rather than to its rheology. Points at which $|\eta^*|$ fell more than 1\% below the highest value already measured at higher frequency were excluded from the Cross fits. This removed one point from FEP at 370~°C, one from HT at 300~°C, and three from HT at 315~°C. Torque at these points was at least 90~$\mu$N$\cdot$m, so instrument resolution was not the cause.

\begin{figure}
    \centering
    \includegraphics[width=1\linewidth]{figures/master_curves_row.pdf}
    \caption{Storage ($G'$, triangles) and loss ($G''$, squares) modulus master curves for FEP 9835 and PFA 440 (reference 330~°C) and HT 2181 (reference 300~°C), assembled from frequency sweeps at three temperatures by horizontal shifting only.}
    \label{fig:mastercurves}
\end{figure}

\begin{figure}
    \centering
    \includegraphics[width=1\linewidth]{figures/cross_viscosity_row.pdf}
    \caption{Shifted viscosity data for each polymer fit to a Cross model.}
    \label{fig:cross_viscosity}
\end{figure}



\begin{table}[htbp]
\caption{Cross model parameters for each polymer, with $a_T$  defined relative to $T_{\mathrm{ref}}$.  Values with $\pm$ are 95\% confidence intervals reported from \texttt{lmfit} \cite{newville_lmfit_2016}.}
\centering
\begin{tabular}{l|ccccc}
\toprule
Polymer & $T_{\mathrm{ref}}$ ($^\circ$C) & $E_a$ (kJ/mol) & $\eta_0$ (Pa$\cdot$s) & $\lambda$ (s) & $n$  \\
\midrule
FEP & 330 & 51.4 & 3044 $\pm$ 41 & 0.0103 $\pm$ 0.0007 & 0.685 $\pm$ 0.029 \\
PFA  & 330 & 44.1 & 4214 $\pm$ 47 & 0.0132 $\pm$ 0.0007 & 0.692 $\pm$ 0.022 \\
HT 2181 & 300 & 92.7 & 12780 $\pm$ 725 & 1.0 $\pm$ 0.2 & 0.529 $\pm$ 0.015 \\
\bottomrule
\end{tabular}
\label{tab:cross_model_parameters}
\end{table}

\section*{Numerical methods}

% NEW
The flow in the die land was treated as isothermal, fully developed annular flow with no slip at the moving wire ($u = V$ at $r = R_w$) and at the stationary die wall ($u = 0$ at $r = R_d$). The axial momentum balance integrates to
\begin{equation}
    \tau_{rz}(r) = -\frac{G r}{2} + \frac{C}{r}, \qquad G = -\frac{dp}{dz},
    \label{eq:stress}
\end{equation}
and since $r\tau_{rz}$ is continuous across the interface, the same $G$ and $C$ hold in both layers. At a given $V$, the solver guesses the two wall stresses and the interface radius, computes $G$ and $C$ from the wall stresses, inverts each layer's Cross model to get $du/dr$ from $\tau_{rz}$, and integrates outward from the wire to get the velocity profile and each layer's flow rate. It then adjusts the three unknowns with trust-region least squares until $u(R_d) = 0$ and each layer's flow rate matches the room-temperature thickness through mass conservation, $Q_k = (\rho_s/\rho_m)_k\,V\pi(r_k^2 - r_{k-1}^2)$. Temperature enters only through the shift factor, $\eta(\dot\gamma, T) = a_T\,\eta_\text{ref}(a_T\dot\gamma)$. Eq.~\ref{eq:stress} also means $|\tau_{rz}|$ has no interior maximum in the gap: any stationary point of $\tau_{rz}$, and the zero-stress point, is a minimum of $|\tau_{rz}|$. So the melt-fracture check only needs the stress at the wire and at the die wall.

For each design, the line speed $V^*$ was defined as the speed at which the larger wall shear stress reached the melt-fracture limit ($2\times10^{5}$~Pa), found with Brent's method on $\ln V$. The fluoropolymer thickness was fixed at its 0.20~mm minimum, and the HT thickness was optimized over the feasible range ($0 < t_\text{HT} \le 0.25 - t_\text{FEP}$~mm) using Brent's method\cite{brent_algorithms_2002} for bounded scalar minimization (as implemented in \texttt{scipy}\cite{2020SciPy-NMeth}). The root-finding step relies on the maximum wall stress increasing monotonically with $V$ (Fig.~\ref{fig:wall_stress_vs_V} (a)), so there is a single crossing of the limit. The bounded minimization relies on $V^*(t_\text{HT})$ having a single peak, which a grid scan confirmed at every temperature (Fig.~\ref{fig:vstar_vs_tHT}).

The provided notebooks were ported to a Python class and modified to accept a Cross model fit instead of a power law.\footnote{Codebase available freely on GitHub at \hyperlink{https://github.com/kylecapybara/667-polymer-project}{https://github.com/kylecapybara/667-polymer-project}} This was motivated by the shear rates in the solution spanning a range of shear rates poorly described by a power law fit (Fig.~\ref{fig:wall_stress_vs_V} (b)).

% NEW
In the single-layer FEP case, the shear rate in the land runs from about 0.08~s$^{-1}$ to 160~s$^{-1}$. Most of that range sits on the Newtonian plateau, and only the top decade reaches the shear-thinning region. Inside the land the shear rate also passes through zero where the velocity peaks, and a power law predicts infinite viscosity there. A power law matched at the wall shear rate overpredicts the FEP viscosity at $\dot\gamma = 1$~s$^{-1}$ by a factor of about 3. The Cross model covers the whole range with one fit, so no shear-rate window has to be chosen.

\begin{figure}
    \centering
    \includegraphics[width=0.49\linewidth]{figures/wall_stress_vs_V.pdf}
    \includegraphics[width=0.49\linewidth]{figures/cross_vs_powerlaw_fep_only.pdf}
    \caption{(a) Wall shear stress magnitude as a function of line speed at the wire (FEP) and at the die wall (HT) for the optimal two-layer design at 330~°C. Both increase monotonically with $V$. (b) Cross model viscosity of FEP at 330~°C. The shaded band marks the range of shear rates FEP experiences in the die land at $V^*$ for the single-layer, 0.20~mm base case.}
    \label{fig:wall_stress_vs_V}
\end{figure}



\section*{Base case: pure fluoropolymer}

Since the wall stress is proportional to the viscosity at the wall, and viscosity decreases with temperature, the highest allowable temperature was selected. Maximum line speeds for the base case of a single pure fluoropolymer layer coating are shared in \cref{tab:basecase}. The base cases point towards FEP allowing faster line speeds before melt fracture, and thinner fluoropolymer layers allowing higher line speeds before melt fracture.

\begin{table}[h]
\caption{Single-layer base cases at 330\,$^\circ$C}
\centering
\begin{tabular}{@{}lcc@{}}
\toprule
Polymer & Thickness (mm) & $V_{\max}$ (m/s) \\
\midrule
Pure FEP 9835 & 0.20 & 0.0165 \\
Pure PFA 440 & 0.20 & 0.0114 \\
Pure FEP 9835 & 0.25 & 0.0109 \\
Pure PFA 440 & 0.25 & 0.0075 \\
\bottomrule
\end{tabular}
\label{tab:basecase}
\end{table}

% NEW
FEP beats PFA at both thicknesses by about 45\%, which follows from its lower zero-shear viscosity at 330~°C (3044 vs.\ 4214~Pa$\cdot$s). The two have nearly the same Cross exponent, so neither thins much more than the other at the wall. A thinner coat means less melt has to be pushed through the same land, so the pressure-driven part of the flow is weaker and the wall stress is lower at a given speed. These results set the starting point for the two-layer design: FEP, at its 0.20~mm minimum, on the wire.


\section*{Optimization}

With the temperature set to 330 °C, the fluoropolymer set to FEP, the layer order set to FEP on the inside, and the thickness of the FEP layer set to the minimum of 0.20 mm, this becomes a 1-D optimization problem. This problem is formalized below as the objective being to maximize the wire velocity over the set of possible HT thickness values where the coating satisfies all the constraints. First, the total thickness must be less than 0.25~mm. Additionally, the draw-down ratio can be a maximum of two, and of course the wall shear stress must be less than the limit for melt fracture. The limit applies both at the wire and the die, so here the constraint is written as the maximum wall shear stress at either wall must be less than the limit ($2\times 10^5$~Pa).

% \begin{itemize}
% \item Discrete: $\Fl\in\{\FEP,\PFA\}$, and the layer order $(\phi_1,\phi_2)\in\{(\Fl,\HT),(\HT,\Fl)\}$, where layer $1$ is the inner layer on the wire.
% \item Continuous: line speed $V_w>0$, die temperature $T$ (K), and room-temperature thicknesses $t_{\Fl}$ and $t_{\HT}$.
% \end{itemize}


\begin{align*}
\max_{V_w,\,t_{\Fl},\,t_{\HT}}\quad & V_w \\
\text{s.t.}\quad
% & 558.15\ \mathrm{K} \le T \le 603.15\ \mathrm{K} && \text{(285--330\,}^\circ\text{C)} \tag{C1}\\
% & t_{\Fl} \ge 0.20\ \mathrm{mm} && \text{fluoropolymer thickness} \tag{C2}\\
& t_{\HT} > 0 && \text{HT layer present} \tag{C1}\\
& t_{\Fl}+t_{\HT} \le 0.25\ \mathrm{mm} && \text{total thickness} \tag{C2}\\
& \dfrac{\pi\,(R_d^2-R_w^2)\,V_w}{Q_1+Q_2} \le 2 && \text{draw-down ratio} \tag{C3}\\
& \max_{\xi_{k-1}\le r\le\xi_k}\,|\tau(r)| \le \tau_{\text{melt fracture}} \quad k=1,2 && \text{melt fracture} \tag{C4}\\
\end{align*}

% NEW
We checked the choice to hold FEP at 0.20~mm by re-optimizing the HT layer at thicker FEP layers (\cref{tab:fep_thickness}). $V^*$ drops each time FEP gets thicker, so the 0.20~mm lower bound is an active constraint and fixing it there loses nothing.

\begin{table}[h]
\caption{Optimal line speed at 330\,$^\circ$C as the FEP layer is made thicker, with the HT thickness re-optimized each time.}
\centering
\begin{tabular}{@{}ccc@{}}
\toprule
$t_\text{FEP}$ (mm) & optimal $t_\text{HT}$ (mm) & $V^*$ (m/s) \\
\midrule
0.20 & 0.029 & 0.0570 \\
0.21 & 0.032 & 0.0506 \\
0.22 & 0.030 & 0.0455 \\
0.23 & 0.020 & 0.0391 \\
\bottomrule
\end{tabular}
\label{tab:fep_thickness}
\end{table}


\section*{Results \& Discussion}

% NEW
Adding an HT layer outside the FEP raises the speed limit at every temperature (Fig.~\ref{fig:vstar_vs_tHT}). At 330~°C, even a 0.002~mm HT skin lifts $V^*$ from 0.0165~m/s for FEP alone to about 0.029~m/s. $V^*$ keeps climbing up to roughly 0.03~mm of HT and then falls off slowly. Two effects compete here. HT is far more shear-thinning than FEP ($\lambda \approx 1$~s against 0.01~s, and a lower exponent), so at the die-wall shear rates it is a thin, low-viscosity layer next to the stationary wall. At the optimum the HT viscosity at the die wall is about 200~Pa$\cdot$s, while the FEP next to it is around 1500~Pa$\cdot$s (Fig.~\ref{fig:profiles_at_optimum}). The HT layer acts as a lubricating skin that carries most of the velocity drop to zero. But every bit of HT added is more melt that has to be pushed through a land that is already full, so the pressure gradient rises, and with it the stress at the die wall. Past about 0.03~mm the second effect wins. At 330~°C the curve drops more steeply above about 0.048~mm, where the limiting wall switches from the die wall to the wire.

$V^*$ increased monotonically with die temperature over 285--330~°C, so the optimum lies at 330~°C, giving $V^* = 0.057$~m/s (3.4~m/min) at $t_\text{HT} = 0.029$~mm. The draw-down ratio (1.15) and the other constraints were verified at the optimum.

% NEW
That is 3.46 times the speed of a 0.20~mm FEP coat on its own at the same temperature (Fig.~\ref{fig:vstar_vs_T}). The advantage grows with temperature, from a factor of 1.96 at 285~°C to 3.46 at 330~°C. This follows from the activation energies in \cref{tab:cross_model_parameters}. HT's viscosity falls about twice as fast with temperature as FEP's, so the lubricating layer gets more effective as the die heats up. The optimal HT thickness barely moves across the sweep (0.025~mm at 285~°C to 0.029~mm at 330~°C), so a single design works over the whole window if the die temperature has to come down.

The draw-down ratio is not a concern for this design. With the land completely filled, it is set by the thicknesses alone and does not depend on $V$, and at 1.15 it is well under the limit of 2. The total coating is 0.229~mm, also under the 0.25~mm limit. The constraint that sets $V^*$ is melt fracture at the die wall, in the HT layer.

% NEW
The radial profiles at the optimum (Fig.~\ref{fig:profiles_at_optimum}) show how the two layers split the work. The shear stress runs from $+1.05\times10^5$~Pa at the wire to $-2.0\times10^5$~Pa at the die wall, so the wire sits at only about half the melt-fracture limit. The stress is positive near the wire because the pressure flow pushes the FEP ahead of the wire. The velocity peaks at about 0.059~m/s near $r = 0.38$~mm, slightly faster than the wire itself. The FEP stays close to its zero-shear viscosity across most of its thickness, with shear rates up to about 100~s$^{-1}$. The HT layer, in contrast, sees 600--1000~s$^{-1}$, where it has thinned by more than an order of magnitude. In the land the FEP layer is 0.275~mm thick and the HT layer 0.062~mm, and the HT layer takes almost all of the drop from 0.05~m/s to zero at the die wall.

% NEW
Several assumptions limit how far these numbers should be trusted. The HT data stop at 315~°C, so the HT viscosity at 330~°C is an Arrhenius extrapolation from three temperatures. If the design must stay within the measured range for every polymer, the 315~°C result applies instead: $V^* = 0.036$~m/s (2.2~m/min) at $t_\text{HT} = 0.028$~mm, still 2.9 times the FEP-only speed. The model assumes fully developed flow, but the land is short ($L/h \approx 0.3$), so entrance effects in the real die could change the wall stresses. The flow is also treated as isothermal, with no viscous heating, die swell, or elastic effects, and the interface between the layers is assumed stable. Interfacial instability is a known problem in coextrusion when the viscosity ratio between layers is as large as it is here (about 6 at the interface). Finally, melt fracture is taken as a sharp cutoff at $2\times10^5$~Pa at either wall, and the same critical stress is applied to both polymers.

\begin{figure}
    \centering
    \begin{subfigure}[b]{0.48\linewidth}
        \centering
        \includegraphics[width=\linewidth]{figures/vstar_vs_tHT.pdf}
        \caption{Varying HT layer's thickness}
        \label{fig:vstar_vs_tHT}
    \end{subfigure}
    \hfill
    \begin{subfigure}[b]{0.48\linewidth}
        \centering
        \includegraphics[width=\linewidth]{figures/vstar_vs_T.pdf}
        \caption{Varying the die temperature}
        \label{fig:vstar_vs_T}
    \end{subfigure}
    \caption{Wire coating speed limits with a 0.20~mm FEP inner layer. (a) $V^*$ against HT thickness at four die temperatures. Each curve has a single peak near 0.03~mm. (b) $V^*$ against die temperature with the HT thickness re-optimized at each point, compared to 0.20~mm of FEP alone.}
    \label{fig:vstar_combined}
\end{figure}

\begin{figure}
    \centering
    \includegraphics[width=1\linewidth]{figures/profiles_at_optimum.pdf}
    \caption{Radial profiles of axial velocity, shear stress, and viscosity in the die land for the optimal coating (330~°C, 0.20~mm FEP, 0.029~mm HT, $V^* = 0.057$~m/s). The wire is at the bottom of each panel and the die wall at the top. Shading marks each layer.}
    \label{fig:profiles_at_optimum}
\end{figure}


% NEW
\section*{Conclusions}

A 0.20~mm FEP 9835 layer on the wire with a 0.029~mm HT 2181 layer outside it, run at a die temperature of 330~°C, allows a line speed of 0.057~m/s (3.4~m/min) before melt fracture. That is about 3.5 times faster than the best single-layer fluoropolymer coating. The gain comes from HT's much stronger shear thinning: placed against the die wall, it takes up the high-shear region and lets the FEP flow at low stress. FEP was preferred over PFA because of its lower viscosity, and the FEP layer was kept at its minimum thickness because any extra FEP lowers the speed limit. The optimum sits at the top of the temperature window, which relies on extrapolating the HT rheology 15~°C past the measured data. If that extrapolation is not acceptable, running at 315~°C still gives 2.2~m/min. Measuring HT at 330~°C and checking interfacial stability between the two layers would be the next steps before trying this design on a line.


\pagebreak
\printbibliography
\pagebreak
\input{appendices.tex}
\end{document}
